\documentclass{article}
\usepackage[a4paper, total={6.5in, 10in}]{geometry}
\setlength{\parindent}{0pt}
\usepackage{tcolorbox}
\tcbset{colback=white!94!black, colframe=black!60!white, coltitle = white, boxrule=0.8pt, arc=2mm}
\newtcolorbox{boxs}[2][]{title= #2,#1}
\usepackage{datetime2}
\usepackage[british]{babel}
\usepackage{amsfonts}
\usepackage{amsmath}

\title{\textbf{Advanced Calculus HW}}
\author{Robert O'Connell}
\date{\today}
\begin{document}
\maketitle
\vspace{5mm}

\subsection*{Question 1}
\subsubsection*{Parameterization}

To solve this problem, it may be useful to first condsider a similar but easier problem, such as a circle of radius 5, lying on 
the \(x\)-\(y\) plane.

The parameterization of this curve would simply be
\[
\vec{r_1}(t) = (5\cos t, 5 \sin t, 0) \quad t \in (0,2\pi)
\]
since sine and cosine have periods of \(2\pi\).
\\

Now to get the circle of radius 5 which is in the plane perpendicular to \((1,1,1)/\sqrt{3}\), we will constuct a linear mapping, which 
rotates our simple circle to the desired circle.

To figure out what our mapping should be, we should figure out where the mapping should take the three standard basis vectors of \(\mathbb{R}^3\).
\\

We know the basis vector (0,0,1) which is perpendicular to the simple circle should stay perpendicular to the circle and thus get mapped to 
\((1,1,1)/\sqrt{3}\). 

To find where the other two basis vectors go, we need to find two unit vectors which lie in the plane of the desired circle which are 
perpendicular to each other and to \((1,1,1)/\sqrt{3}\).

One such unit vector \(\vec{v_1}\) could be attained using the dot product
\[
(1,1,1)/\sqrt{3} \ \cdot \ \vec{v_1} = 0
\]
\[
v_{1x} + v_{1y} + v_{1z} = 0
\]
(infinite such vectors exist)
\[
\vec{v_1} = (2,-1,-1)/\sqrt{6}
\]

To find the final vector we use the cross product, since we need a vector perpendicular to the two other vectors
\begin{align*}
    \vec{v_2} &= ((2,-1,-1)/\sqrt{6}) \times ((1,1,1)/\sqrt{3}) \\
    &= \frac{1}{3\sqrt{2}}(\det \begin{bmatrix}
        \hat{i} & \hat{j} & \hat{k} \\
        2 & -1 & -1 \\
        1 & 1 & 1 
    \end{bmatrix}) \\
    &= (0, -3, 3)/3\sqrt{2} \\
    &= (0,-1,1)/\sqrt{2}
\end{align*}
\\

We now know where this linear mapping takes the three basis vectors, so we can display this in a matrix as a linear transformation, with the images
of the standard basis vectors being the columns of the matrix

\[
T: \mathbb{R}^3 \rightarrow \mathbb{R}^3
\]
\[
T(\vec{x}) = A\vec{x}
\]
\[
A = \begin{bmatrix}
    \frac{2}{\sqrt{6}} & 0 & \frac{1}{\sqrt{3}}   \\
    -\frac{1}{\sqrt{6}} & -\frac{1}{\sqrt{2}} & \frac{1}{\sqrt{3}} \\
    -\frac{1}{\sqrt{6}} & \frac{1}{\sqrt{2}} & \frac{1}{\sqrt{3}} 
\end{bmatrix}
\]

\newpage
As a final step, all we need to do is multiply our simple parameterization
by our transformation matrix.
\[
\vec{r_2}(t) = A \begin{bmatrix}
    5\cos t \\
    5 \sin t \\
    0
\end{bmatrix}
\]

\[
\vec{r_2}(t) = \begin{bmatrix}
    \frac{10\cos t}{\sqrt{6}}  \\
    -\frac{5\cos t}{\sqrt{6}} - \frac{5 \sin t}{\sqrt{2}} \\
    -\frac{5\cos t}{\sqrt{6}} + \frac{5 \sin t}{\sqrt{2}}
\end{bmatrix} \quad t \in (0,2\pi)
\]
Which is our final parameterization of the desired circle.

\subsubsection*{Opposite Orientation}

To find a parameterization of opposite orientation we need to find a bijective function which maps \(t \rightarrow \tau\) 
such that \(t = \gamma(\tau)\), for the change in parameter to change the orientation we need \(\gamma '(\tau) < 0\)
\\

One such function is \(\gamma(\tau) = - \tau\) with \(\tau \in (-2\pi , 0)\)

Recall, 
\[
\cos (-x) = \cos (x), \quad \sin(-x) = -\sin(x)
\]

Our new parameterization is,
\[
\vec{r_3}(\tau) = \begin{bmatrix}
    \frac{10\cos \tau}{\sqrt{6}}  \\
    -\frac{5\cos \tau}{\sqrt{6}} + \frac{5 \sin \tau}{\sqrt{2}} \\
    -\frac{5\cos \tau}{\sqrt{6}} - \frac{5 \sin \tau}{\sqrt{2}}
\end{bmatrix} \quad \tau \in (-2\pi,0)
\]

\subsubsection*{Shifting the Circle}

To shift the circle by a given amount in each direction say \((s_x,s_y,s_z)\), simply add each component to the parameterization
to get
\[
\vec{r_4}(t) = \begin{bmatrix}
    \frac{10\cos t}{\sqrt{6}} + s_x \\
    -\frac{5\cos t}{\sqrt{6}} - \frac{5 \sin t}{\sqrt{2}} + s_y \\
    -\frac{5\cos t}{\sqrt{6}} + \frac{5 \sin t}{\sqrt{2}} + s_z
\end{bmatrix} \quad t \in (0,2\pi)
\]


\subsection*{Question 2}
\[
\vec{r}(t) = (2\cos(t^2) -1,2\sin(t^2), 3)
\]

\begin{itemize}
    \item We know that both sine and cosine have periods of \(2\pi\), thus the curve starts to repeat
    after \(t^2\) increases by \(2\pi\), thus we need
    \(t^2 = 2\pi \Rightarrow t = \sqrt{2\pi}\) thus \(a = \sqrt{2\pi}\)
    
    \item To find the derivative of the vector valued function we simply differentiate it component-wise
    \[
    \vec{r} \ ^{\prime}(t) = (-4t\sin(t^2), 4t\cos(t^2), 0)
    \]
    The norm is then
    \begin{align*}
        |\vec{r} \ ^{\prime}(t)| &= \sqrt{(-4t\sin(t^2))^2 + (4t\cos(t^2))^2} \\
        &= \sqrt{16t^2(\cos^2(t^2) + \sin^2(t^2))} \\
        &= 4|t|
    \end{align*}

    \item If \(t = \gamma(\tau) = \sqrt{\tau}\), then \(\vec{r}(\tau) = (2\cos(\tau) -1,2\sin(\tau), 3)\) with \(\tau \in (0,2\pi)\)
    \[
    \frac{d\vec{r}}{d\tau} = \frac{d\vec{r}}{dt} \frac{dt}{d\tau}
    \]
    Where \(\frac{d\vec{r}}{dt} = \vec{r} \ ' (\sqrt{\tau})\) and
    \[
    t = \sqrt{\tau} \Rightarrow \frac{dt}{d\tau} = \frac{1}{2\sqrt{\tau}}
    \]
    Thus
    \[
    \frac{d\vec{r}}{d\tau} = \frac{1}{2\sqrt{\tau}} \vec{r} \ '(\sqrt{\tau})
    \]
    It's norm is then 
    \begin{align*}
        \left|\frac{d\vec{r}}{d\tau}\right| &= \left|\frac{1}{2\sqrt{\tau}} \vec{r} \ '(\sqrt{\tau})\right| \\
        &= \frac{1}{2\sqrt{\tau}}(4\sqrt{\tau}) \\
        &= 2
    \end{align*}
\end{itemize}

\subsection*{Question 3}
\[
\vec{r}(t) = \begin{cases}
    (a,b,c) \quad \text{if } t = 3 \\
    (4\sin(\ln(t+1)), -\sqrt{16-t^2}, \sinh(3t)) \quad \text{otherwise}
\end{cases}
\]

\begin{itemize}
    \item The natural domain of a vector valued function is the intersection of the
    natural domains of the component functions. Remember sine and sinh can take any input, the
    natural log requires inputs greater than zero and square roots can not take in negatives
    \[
    t+1 > 0, \quad 16-t^2 \geq 0, \quad t \in \mathbb{R}
    \]
    \[
    t > -1, \quad |t| \leq 4, \quad t \in \mathbb{R}
    \]
    Thus the natural domain of \(\vec{r}(t)\) is \((-1, 4]\)   
    
    \item If a function is continuous at a point \(a\) then,
    \[
    \lim_{x\to a}{f(x)} = f(a)
    \]
    Also, for a vector valued function then each of the limits of the components must 
    be equal.
    Since for \(\vec{r}(t)\) each of the components themselves are continuous then we know
    \[
    a = x(3), \quad b = y(3), \quad c = z(3)
    \]
    \[
    a = 4 \sin(\ln(4)), \quad b = -\sqrt{7}, \quad c = \sinh(9)
    \]

    \item We know in general that differentability implies continuity while the converse of the statement is
    not neccesarily true,
    but in this case, \(\vec{r}(t)\) is differentiable, since all of it's component functions are
    differentiable at \(t = 3\), the fact that \(\vec{r}(t)\) is defined seperately at \(t=3\) has no effect, since the curve
    remains continuous.

    \[
    \vec{r} \ '(t) = \left(\frac{4}{t+1}\cos(\ln(t+1)), \frac{t}{\sqrt{16-t^2}}, 3\cosh(3t)\right) 
    \]
\end{itemize}
\end{document}